\begin{table}[tbp]\centering
\caption{How Often the Four Readings Give Different Judgments}\label{tab:four}
\resizebox{\ifdim\width>\textwidth\textwidth\else\width\fi}{!}{%
\begin{tabular}{lrrr}
\toprule
 & Percent of months & 90\% interval & Months \\
\midrule
All four readings have the same sign & 81.6 & [70.2, 82.9] & 301 \\
At least two materially opposing readings, $t$ = 10 bp & 10.3 & [8.4, 19.8] & 38 \\
At least two materially opposing readings, $t$ = 25 bp & 2.2 & [1.6, 8.1] & 8 \\
Four-reading range at least 10 bp & 92.4 & [87.5, 97.6] & 341 \\
Four-reading range at least 25 bp & 70.5 & [58.5, 88.9] & 260 \\
Market (futures factor) and FRB/US materially opposing, $t$ = 25 bp & 2.2 & [0.0, 5.4] & 8 \\
\bottomrule
\end{tabular}}
\begin{flushleft}\footnotesize\textit{Notes:} Monthly 12-month changes in the three assets relative to $r^*$, January 1996 to September 2026 (369 months, overlapping), read with the four portfolios on the common two-year-yield scale. Materially opposing: at least one reading at least $t$ above zero and at least one at least $t$ below zero. Four-reading range: the largest less the smallest of the four readings. Intervals: 5th and 95th percentiles over 5,000 bootstrap draws of the announcements (market weights and the common response vector redrawn; FRB/US and GS weights fixed), conditional on the observed history.\end{flushleft}
\end{table}
